\section{Gr\"obner, Rees and marked-channel structures}\label{sec:gf1}

The dark ideal \(I=(L,Q)\) records exact zero sets but forgets which generator
enters at which perturbation and time depth.  The following hierarchy makes the
information loss explicit.

Introduce adapted coordinates
\[
\lambda=L,\qquad \xi=d_2-d_3,\qquad \eta=d_4-d_3,\qquad \zeta=d_3.
\]
Then
\[
d_1=\lambda+\xi+\eta+\zeta,\quad d_2=\xi+\zeta,
\quad d_3=\zeta,\quad d_4=\eta+\zeta
\]
and
\begin{equation}\label{eq:normalcrossing}
Q=\xi\eta-\lambda\zeta,\qquad I=(\lambda,\xi\eta)
=(\lambda,\xi)\cap(\lambda,\eta).
\end{equation}
Thus the two components meet normally inside the hyperplane \(\lambda=0\), but
their union has codimension two in the ambient four-space; calling it an
ambient normal-crossing divisor would be incorrect.

\begin{theorem}[Gr\"obner fan of the dark ideal]\label{thm:fan}
In the original coordinates the dark ideal has 16 full-dimensional monomial
initial cones.  Twelve initial ideals are radical squarefree intersections of
two planes.  Four are doubled planes:
\[
(d_1,d_3^2),\quad(d_2,d_4^2),\quad(d_3,d_1^2),\quad(d_4,d_2^2).
\]
Every cone has dimension \(2\), degree \(2\), Hilbert series
\((1+t)/(1-t)^2\) and Hilbert polynomial \(2n+1\).  Initialisation commutes with
the two-component intersection in exactly the twelve radical cones.
\end{theorem}
\begin{proof}
Run Buchberger reduction for each strict ordering chamber of the four variables,
merge chambers with the same reduced marked basis, and retain the 16 distinct
initial ideals.  Monomial primary decomposition gives the twelve squarefree and
four primary cases.  Counting standard monomials yields the common Hilbert
series and polynomial.  The complete marked bases, representative weights,
facets and adjacency graph are reproduced in \cref{app:atlases}.  An independent
exact script checks the counts and invariants against the released JSON atlas.
\end{proof}

\begin{proposition}[Rees algebra]\label{prop:rees}
Let \(Y_L,Y_Q\) mark the two generators.  Then
\[
\mathcal R(I)\cong
\frac{A[Y_L,Y_Q]}{(QY_L-LY_Q)}.
\]
The ideal is of linear type and its powers are integrally closed.  The blow-up
does not separate the two dark components: its exceptional geometry retains a
singular section over their intersection.
\end{proposition}
\begin{proof}
The two generators form a regular sequence in the polynomial ring, so the Rees
kernel is generated by the Koszul relation \(QY_L-LY_Q\).  In the adapted
normal form \((\lambda,\xi\eta)\), the Newton polyhedron of each monomial power
is saturated in its exponent lattice, giving integral closure.  The two affine
blow-up charts follow by setting either Rees coordinate to one; their equations
retain the component-intersection section.
\end{proof}

For an arc with orders
\(a=\ord\lambda\), \(b=\ord\xi\), \(c=\ord\eta\), its basic contact profile is
\begin{equation}\label{eq:contact}
\bigl(\min(a,b+c),\min(a,c),\min(a,b),\min(a,b,c)\bigr).
\end{equation}
Arc-space contact loci provide the general geometric context
\parencite{Ein2004}; \cref{eq:contact} is the exact specialised formula.

\subsection{Tagged ideal and marked section}

Introduce tags \(u\) for perturbation order and \(x\) for time depth.  The
tagged channel ideal is
\begin{equation}
J=(Lu x^3,\;Qu^2x^4,\;Ru^3x^5)=(Lu x^3,Qu^2x^4),
\end{equation}
with first syzygy \((Qux,-L)\).  Saturating away the tags recovers the dark
ideal, but the plain ideal cannot reconstruct the tags.

Let a free module have basis \(e_L,e_Q,e_R\).  The numerator marking is
\[
s_N=(1+2x)e_L+2(1+x)e_Q+e_R,
\]
and the kernel includes
\[
(Qux,-L,0),\qquad(d_2d_4u^2x^2,(d_2+d_4)ux,1).
\]
The first records the two-generator relation with depth tags; the second records
the cubic syzygy.  Even this marked numerator section is not the full amplitude:
denominator inversion and factorial phases remain necessary.

\begin{figure}[ht]
\centering\begin{tikzpicture}[node distance=7mm,>=Latex,
  box/.style={draw,rounded corners,fill=blue!7,minimum width=43mm,minimum height=8mm,align=center,font=\small},
  arr/.style={->,thick,red!65!black}]
\node[box,fill=green!10] (amp) {full amplitude section\\phases and denominator};
\node[box,below=of amp] (marked) {marked numerator module\\generator provenance};
\node[box,below=of marked] (tag) {tagged ideal\\perturbation/time depths};
\node[box,below=of tag] (rees) {Rees algebra\\approach direction};
\node[box,below=of rees] (plain) {plain dark ideal\\zero locus};
\draw[arr] (amp)--node[right,font=\scriptsize]{forget} (marked);
\draw[arr] (marked)--node[right,font=\scriptsize]{forget} (tag);
\draw[arr] (tag)--node[right,font=\scriptsize]{forget tags} (rees);
\draw[arr] (rees)--node[right,font=\scriptsize]{project} (plain);
\end{tikzpicture}

\caption{Information hierarchy.  Moving left is a lossy forgetful operation;
the plain dark ideal cannot recover depth tags or amplitude phases.}
\label{fig:hierarchy}
\end{figure}

For the mixed path \(d=(s,s^3,-s,s^4)\), a tied face has
\[
W_M=-c^3+2c^4,\qquad W_A=-\frac{i c^3}{6}+\frac{c^4}{12}.
\]
At \(c=2i\), the amplitude face cancels and the exact next coefficient is
\(-32/15\) at order seven.  This control is \evidence{VERIFIED-EXACT} and
exhibits why the full amplitude section is the terminal object in the hierarchy.
